\chapter{Régulateurs et agréments}\label{ch:fm:regulators-and-licences}

Une entreprise qui ne négocie que son propre argent n'a pas à être cliente de qui que ce soit, et elle n'a pas de clients. Savoir si elle a besoin d'un \omterm{def:fm:regulators-and-licences:perimeter}{agrément} est une
autre affaire. La deuxième directive sur les marchés d'instruments financiers conserve une exemption pour les entreprises qui négocient pour compte propre dans l'Union européenne,
puis la retire à presque toutes les entreprises que cette série décrit~: aux teneurs de marché, aux membres d'une plateforme de négociation et aux entreprises ayant un accès électronique direct
à l'une d'elles, et à toute entreprise qui applique une technique de trading algorithmique à haute fréquence, quelle que soit sa taille et bien qu'elle
n'ait aucun client. Les lieux où une entreprise peut négocier, et ce qu'elle doit être pour le faire, sont fixés par des règles qui varient selon le pays et selon l'activité, et la carte de
ces règles fait partie de la conception de l'entreprise.

\section{Pourquoi une société de trading a besoin d'un agrément}

La réglementation suit les activités, et non les entreprises. Une même entreprise peut ne pas être réglementée pour une de ses activités, être enregistrée pour une autre et agréée pour une troisième,
et un nouveau desk, un nouveau produit ou un nouveau bureau peut la faire franchir une ligne qu'elle ignorait exister.

\begin{definition}[Périmètre réglementaire, agrément]\label{def:fm:regulators-and-licences:perimeter}
Le \emph{périmètre réglementaire}\index{périmètre réglementaire} d'une juridiction est la frontière entre les activités qui requièrent l'agrément
d'un régulateur et celles qui n'en requièrent pas, tracée par les définitions des activités réglementées de la législation et par ses exemptions. L'\emph{agrément}\index{agrément}
est la permission, accordée par un régulateur sur demande, d'exercer des activités réglementées déterminées, avec les conditions, les \omterm{def:fm:the-proprietary-market-making-firm:ownfunds}{fonds propres} et
la surveillance qui s'y rattachent.
\end{definition}

\begin{definition}[Négociation pour compte propre, entreprise d'investissement]\label{def:fm:regulators-and-licences:own}
La \emph{négociation pour compte propre}\index{négociation pour compte propre} consiste à négocier avec des capitaux propres, ce qui donne lieu à des
transactions sur des instruments financiers. Une \emph{entreprise d'investissement}\index{entreprise d'investissement} est, dans le droit européen, une personne morale dont l'activité habituelle ou professionnelle consiste à fournir des services d'investissement à des tiers
ou à exercer des activités d'investissement à titre professionnel~; la négociation pour compte propre est l'une de ces
activités.
\end{definition}

L'outil de travail de ce chapitre est une carte de ces règles sous forme de données, chaque ligne désignant une juridiction, une activité, le régime concerné, la référence juridique,
le délai légal de décision lorsqu'il existe, le mois de vérification et les lignes du registre qui en fournissent la source
(\cref{lst:fm:regulators-and-licences:map}). La carte indique où se trouvent les règles~; savoir si une entreprise donnée est dans un périmètre est une
question pour ses avocats et pour son régulateur.

\omcode{code/firm/regmap/firm_regmap.py}{39}{64}{Vérification du schéma de la carte réglementaire, et requête sur le profil d'activités et de
juridictions d'une entreprise.}\label{lst:fm:regulators-and-licences:map}

\section{États-Unis~: la SEC, la CFTC et les organismes d'autorégulation}

Les États-Unis répartissent le champ par instrument. Les titres relèvent de la Securities and Exchange Commission, les contrats à terme, les options sur contrats à terme et les swaps
de la Commodity Futures Trading Commission~; chacune délègue une grande part de l'enregistrement et de la surveillance courante à un organisme de membres.

\begin{definition}[Organisme d'autorégulation]\label{def:fm:regulators-and-licences:sro}
Un \emph{organisme d'autorégulation}\index{organisme d'autorégulation} est un organisme de membres, tel qu'une association de valeurs mobilières ou une bourse,
qui élabore et fait respecter des règles pour ses membres sous la supervision d'un régulateur public, et auquel les entreprises enregistrées sont tenues
d'adhérer.
\end{definition}

\begin{definition}[Swap dealer, commodity pool operator, introducing broker]\label{def:fm:regulators-and-licences:cftc}
Un \emph{swap dealer}\index{swap dealer} se présente comme un dealer en swaps, anime un marché sur ceux-ci, ou conclut régulièrement des swaps avec
des contreparties pour son propre compte dans le cours normal de son activité. Un \emph{commodity pool operator}\index{commodity pool operator} gère un pool
qui réunit les fonds de plusieurs personnes pour négocier des contrats à terme, des options sur contrats à terme ou des swaps, et sollicite des fonds pour ce pool. Un \emph{introducing
broker}\index{introducing broker} sollicite ou accepte des ordres sur des contrats à terme, des options sur contrats à terme ou des swaps, mais n'accepte pas d'argent ni d'autres actifs de la part
de clients pour les soutenir.
\end{definition}

\begin{dated}{2026-09}{Enregistrement aux États-Unis}\label{dat:fm:regulators-and-licences:us}
\textbf{Titres.} Un broker ou un dealer doit s'enregistrer auprès de la SEC pour effectuer des transactions sur titres (Exchange Act, s.~15(a))~; dans les 45 jours
suivant une demande, la SEC accorde l'enregistrement ou engage une procédure, qui doit s'achever dans les 120 jours, prolongeables de 90 au plus (s.~15(b)(1)). Un broker-dealer enregistré doit en outre adhérer
à une association de valeurs mobilières enregistrée, sauf s'il ne négocie que sur les bourses dont il est membre (s.~15(b)(8))~; la Rule 15b9-1, telle que modifiée en septembre 2023,
limite cette exemption aux membres de bourse sans comptes clients qui ne négocient que
sur leurs bourses, avec quelques exceptions étroites. Un conseiller en investissement doit s'enregistrer au titre de l'Advisers Act (s.~203(a)), selon le même calendrier de 45 jours. \textbf{Dérivés.}
Le Commodity Exchange Act impose à certaines entreprises de s'enregistrer auprès de la CFTC, qui a délégué l'enregistrement à la National
Futures Association~; à peu d'exceptions près, les entreprises enregistrées doivent être membres de la NFA. Les catégories comprennent les futures commission merchants, les
\omterm{def:fm:regulators-and-licences:cftc}{introducing brokers}, les \omterm{def:fm:regulators-and-licences:cftc}{commodity pool operators}, les commodity trading advisors et les \omterm{def:fm:regulators-and-licences:cftc}{swap dealers}.
\end{dated}

\begin{omfigure}
\begin{tikzpicture}[font=\footnotesize,>=stealth,box/.style={draw,rounded corners,align=center,minimum height=0.75cm,text width=3.2cm}]
\node[box,fill=gray!15] (sec) at (-3.4,1.6) {SEC\\{\scriptsize titres}};
\node[box,fill=gray!15] (cftc) at (3.4,1.6) {CFTC\\{\scriptsize contrats à terme, options sur contrats à terme, swaps}};
\node[box,fill=omDef!15] (finra) at (-3.4,0) {association de valeurs mobilières\\{\scriptsize règles d'adhésion}};
\node[box,fill=omDef!15] (nfa) at (3.4,0) {NFA\\{\scriptsize enregistrement délégué}};
\node[box,fill=omProp!15] (bd) at (-3.4,-1.6) {broker-dealer, conseiller en investissement};
\node[box,fill=omProp!15] (fcm) at (3.4,-1.6) {FCM, IB, CPO, CTA, swap dealer};
\draw[->] (sec) -- node[left,font=\scriptsize] {supervise} (finra); \draw[->] (cftc) -- node[right,font=\scriptsize] {supervise} (nfa);
\draw[->] (sec.south west) to[bend right=30] node[left,font=\scriptsize] {enregistre} (bd.north west);
\draw[->] (finra) -- node[left,font=\scriptsize] {membres} (bd); \draw[->] (nfa) -- node[right,font=\scriptsize] {enregistre, membres} (fcm);
\end{tikzpicture}

\omcaption{La répartition américaine du champ~: la SEC pour les titres et la CFTC pour les dérivés, chacune avec un organisme de membres auquel
doivent adhérer les entreprises enregistrées. Schéma des règles de l'encadré daté.}\label{fig:fm:regulators-and-licences:us}
\end{omfigure}

Pour une société de trading pour compte propre, la modification de 2023 constitue le point d'accroche à l'américaine~: une entreprise qui négocie hors bourse, même pour elle-même, doit désormais
adhérer à l'association de valeurs mobilières et respecter son règlement de marché. L'accès au marché (Rule 15c3-5, livre~11) impose des contrôles avant transmission, quel que soit celui qui
parraine les ordres de l'entreprise.

\section{Europe et Royaume-Uni}

\begin{definition}[Passeport européen]\label{def:fm:regulators-and-licences:passport}
Le \emph{passporting}\index{passporting} est le droit, pour une entreprise agréée dans un État membre de l'Union européenne ou de l'Espace économique européen, de
fournir dans les autres les services couverts par son \omterm{def:fm:regulators-and-licences:perimeter}{agrément}, sur notification à son régulateur d'origine, sans nouvel \omterm{def:fm:regulators-and-licences:perimeter}{agrément}.
\end{definition}

\begin{dated}{2026-09}{Agrément européen au titre de MiFID II}\label{dat:fm:regulators-and-licences:eu}
La directive 2014/65/UE exige un \omterm{def:fm:regulators-and-licences:perimeter}{agrément} préalable pour les services et activités d'investissement exercés à titre professionnel (art.~5), et le
régulateur doit statuer dans les six mois suivant une demande complète (art.~7(3))~; une entreprise agréée peut fournir ses services dans d'autres États membres
(art.~34). Les entreprises qui négocient uniquement pour compte propre sont exemptées (art.~2(1)(d)), sauf si elles sont teneurs de marché, membres ou participants d'un marché réglementé
ou d'un MTF, si elles ont un accès électronique direct à une plateforme de négociation, si elles appliquent une technique de trading algorithmique à haute fréquence, ou si elles négocient pour compte propre
lorsqu'elles exécutent des ordres de clients. Une entreprise qui pratique le trading algorithmique doit disposer de systèmes et de contrôles des risques efficaces, d'une continuité d'activité et de systèmes testés
(art.~17(1)), et doit notifier aux autorités de son État d'origine et des plateformes où elle négocie (art.~17(2)).
\end{dated}

\begin{dated}{2026-09}{Autorisation au Royaume-Uni}\label{dat:fm:regulators-and-licences:uk}
En vertu du Financial Services and Markets Act 2000, nul ne peut exercer une activité réglementée au Royaume-Uni sans être agréé ou exempté
(s.~19, l'interdiction générale). Une demande d'\omterm{def:fm:regulators-and-licences:perimeter}{agrément} doit être tranchée dans les six mois suivant la réception par le régulateur d'un dossier complet, et un dossier incomplet
dans les douze mois (s.~55V). Les exemptions du Royaume-Uni pour la \omterm{def:fm:regulators-and-licences:own}{négociation pour compte propre} sont les siennes et ne sont pas reprises ici~: la carte signale
les lignes britanniques comme devant être vérifiées.
\end{dated}

La dérogation européenne est rédigée de telle sorte que presque aucune société de trading électronique ne la remplit. Une société pour compte propre qui est membre d'une bourse, ou qui
atteint une plateforme par un accès électronique direct, est exclue~; une société qui fait de la colocalisation, qui génère ses ordres sans intervention humaine et qui envoie de nombreux
messages applique une technique à haute fréquence (la définition du livre~10, reprise dans les termes de la directive, article~4(1)(40)) et est de nouveau exclue. Ce qu'elle y gagne
en contrepartie, c'est le passeport européen.

\begin{omfigure}
\begin{tikzpicture}[font=\footnotesize,>=stealth]
\node[draw,rounded corners,fill=gray!10,align=center,text width=4.6cm,minimum height=0.8cm] (q0) at (0,0) {négocie uniquement pour compte propre~?};
\node[draw,rounded corners,fill=omDef!15,align=center,text width=4.6cm,minimum height=0.8cm] (q1) at (0,-1.9) {teneur de marché, membre de bourse, accès électronique direct, technique HFT ou ordres de clients~?};
\node[draw,rounded corners,fill=omThm!20,align=center,text width=2.6cm,minimum height=0.8cm] (ex) at (5.0,-1.9) {exemptée (art.~2(1)(d))};
\node[draw,rounded corners,fill=omProp!20,align=center,text width=4.6cm,minimum height=0.8cm] (au) at (0,-3.9) {agrément comme entreprise d'investissement (art.~5, 7)~: décision dans les six mois};
\node[draw,rounded corners,fill=omProp!10,align=center,text width=2.6cm,minimum height=0.8cm] (pp) at (5.0,-3.9) {passeport vers les autres États membres (art.~34)};
\node[draw,rounded corners,fill=omProp!20,align=center,text width=2.6cm,minimum height=0.8cm] (au2) at (-5.0,-1.9) {agrément pour ses autres services};
\draw[->] (q0) -- node[right,font=\scriptsize] {oui} (q1); \draw[->] (q0) -| node[above,pos=0.2,font=\scriptsize] {non} (au2);
\draw[->] (q1) -- node[above,font=\scriptsize] {non} (ex); \draw[->] (q1) -- node[right,font=\scriptsize] {oui} (au); \draw[->] (au) -- (pp);
\end{tikzpicture}

\omcaption{La dérogation de MiFID II pour la \omterm{def:fm:regulators-and-licences:own}{négociation pour compte propre} d'une entreprise qui négocie des instruments financiers autres que les dérivés de matières premières, telle que la directive
la formule. Presque chaque société de trading électronique répond oui à la deuxième question. Schéma du texte juridique, non un conseil.}\label{fig:fm:regulators-and-licences:tree}
\end{omfigure}

\section{Asie-Pacifique}

Les marchés asiatiques sont régulés pays par pays. À Hong Kong, une société qui exerce une activité réglementée a besoin d'une licence
de la Securities and Futures Commission~; la Securities and Futures Ordinance énumère treize types d'activité, parmi lesquels la négociation de titres (type~1),
la négociation de contrats à terme (type~2), la fourniture de services de trading automatisé (type~7) et la gestion d'actifs (type~9). Savoir si une entreprise qui ne négocie que
propre argent exerce une telle activité dépend des définitions et des exemptions de l'ordonnance, si bien que la carte la signale comme un point à vérifier, et non comme une obligation.
Singapour, le Japon et l'Australie ont chacun leurs propres licences~; le cadre de ce chapitre ne reprend que les lignes pour lesquelles il dispose d'une source, et le dit.

\section{Règles du trading algorithmique}

Une fois dans le périmètre, une entreprise électronique est soumise à des obligations propres aux algorithmes. Ce sont les contrôles du chapitre du livre~11 consacré aux contrôles de risque
inscrits dans la loi~: dans l'Union européenne, l'article~17 de MiFID II et ses normes techniques~; aux États-Unis, la Market Access
Rule pour le courtier dont l'entreprise utilise la connexion (\cref{tab:fm:regulators-and-licences:algo}).

\begin{table}[ht]
\centering\small
\begin{tabular}{>{\raggedright\arraybackslash}p{1.2cm}>{\raggedright\arraybackslash}p{5.4cm}>{\raggedright\arraybackslash}p{6.2cm}}
\toprule
où & obligation & référence \\
\midrule
UE & systèmes et contrôles de risque, limites, pas d'ordres erronés & Directive 2014/65/UE art.~17(1) \\
UE & continuité d'activité~; systèmes testés et surveillés & Directive 2014/65/UE art.~17(1) \\
UE & notification aux autorités d'origine et des plateformes & Directive 2014/65/UE art.~17(2) \\
UE & exigences organisationnelles détaillées & Règlement délégué (UE) 2017/589 (livre~11) \\
US & contrôles de risque avant négociation pour l'accès au marché & Rule 15c3-5 (livre~11) \\
\bottomrule
\end{tabular}
\caption{La liste de contrôle du trading algorithmique de \texttt{firm.regmap.ALGO\_CHECKLIST}, chaque point avec sa référence.}\label{tab:fm:regulators-and-licences:algo}
\end{table}

\section{Tutoriel~: trois bureaux}\label{tut:fm:regulators-and-licences:tut}

\begin{tutorial}
\textbf{Objectif.} Interroger la carte pour une société pour compte propre à haute fréquence qui s'installe à Londres puis ouvre des bureaux à Chicago et à Hong Kong.
\textbf{État final~:} les régimes que touche chaque étape (\cref{tab:fm:regulators-and-licences:profiles}) et les délais légaux de décision sur le parcours
(\cref{fig:fm:regulators-and-licences:periods}).
\begin{tutsteps}
\item \textbf{La carte.} \texttt{fm\_regmap.rows()} charge les 20 lignes de \texttt{data/desk/regmap.csv} (quatre juridictions, dix régimes) et
\texttt{firm.regmap.check} vérifie que chaque ligne a une référence, un mois de mise à jour et une source.
\item \textbf{Les profils.} Londres~: \omterm{def:fm:regulators-and-licences:own}{négociation pour compte propre} sur bourse avec une technique à haute fréquence. Chicago ajoute la même chose sur les titres américains.
Hong Kong ajoute la \omterm{def:fm:regulators-and-licences:own}{négociation pour compte propre} sur bourse.
\item \textbf{Les requêtes.} \texttt{fm\_regmap.profiles()} énumère les lignes, les régimes distincts et le délai légal de décision le plus long.
\item \textbf{La maintenance.} \texttt{firm.regmap.stale(rows, today, 12)} énumère les lignes non vérifiées depuis un an.
\end{tutsteps}
\end{tutorial}

\begin{table}[ht]
\centering\small
\begin{tabular}{lrrrrl}
\toprule
profil & lignes & régimes & requis & à vérifier & décision la plus longue \\
\midrule
Londres & 2 & 1 & 0 & 2 & 6 mois (RU) \\
Londres et Chicago & 5 & 4 & 2 & 2 & 6 mois (RU) \\
Londres, Chicago et Hong Kong & 6 & 5 & 2 & 3 & 6 mois (RU) \\
\bottomrule
\end{tabular}
\caption{La carte réglementaire interrogée pour une société qui ouvre des bureaux (l'étape de Chicago ajoute aussi une obligation, la Market Access Rule). Information,
non conseil. Données~: \texttt{fm\_regmap.profiles}.}\label{tab:fm:regulators-and-licences:profiles}
\end{table}

\begin{omfigure}
\begin{tikzpicture}
\begin{axis}[width=10.5cm,height=4.6cm,xbar stacked,bar width=12pt,xmin=0,xmax=7,xlabel={\small lignes de la carte touchées},ytick={0,1,2},
  yticklabels={Londres,{Londres et Chicago},{Londres, Chicago et Hong Kong}},y dir=reverse,ymin=-0.6,ymax=2.6,
  tick label style={font=\footnotesize},grid=major,grid style={gray!20},table/col sep=comma,
  legend cell align=left,legend style={font=\scriptsize,at={(0.5,-0.35)},anchor=north,/tikz/every even column/.append style={column sep=8pt}},legend columns=3]
\addplot[fill=omProp!55,draw=omProp] table[y=pos,x=required]{figdata/desk/17-regulators-and-licences/profiles.csv};
\addplot[fill=omThm!50,draw=omThm] table[y=pos,x=obligation]{figdata/desk/17-regulators-and-licences/profiles.csv};
\addplot[fill=gray!40,draw=gray] table[y=pos,x=check]{figdata/desk/17-regulators-and-licences/profiles.csv};
\legend{requis,obligation,à vérifier}
\end{axis}
\end{tikzpicture}

\omcaption{Lignes de la carte réglementaire touchées par une entreprise à mesure qu'elle ouvre des bureaux, selon leur statut. Chaque bureau ajoute des régimes, et chaque ligne à vérifier est une question pour
les juristes de l'entreprise. Données~: \texttt{fm\_regmap.profiles}.}\label{fig:fm:regulators-and-licences:profiles}
\end{omfigure}

\begin{omfigure}
\begin{tikzpicture}
\begin{axis}[width=10.5cm,height=5cm,xbar,bar width=9pt,xmin=0,xmax=400,xlabel={\small jours},ytick={0,1,2,3,4,5},
  yticklabels={RU~: dossier complet,RU~: dossier incomplet,UE~: dossier complet,US~: décision ou procédure ouverte,US~: procédure close,US~: avec prolongation},
  y dir=reverse,ymin=-0.6,ymax=5.6,tick label style={font=\footnotesize},grid=major,grid style={gray!20},table/col sep=comma,
  nodes near coords,nodes near coords style={font=\scriptsize}]
\addplot[fill=omDef!55,draw=omDef] table[y=pos,x=days]{figdata/desk/17-regulators-and-licences/periods.csv};
\end{axis}
\end{tikzpicture}

\omcaption{Délais légaux maximaux de décision d'un régulateur sur une demande, en jours (six mois comptés pour 183)~: \omterm{def:fm:regulators-and-licences:perimeter}{agrément} britannique, \omterm{def:fm:regulators-and-licences:perimeter}{agrément} d'\omterm{def:fm:regulators-and-licences:own}{entreprise d'investissement} de l'UE,
et enregistrement d'un broker-dealer aux États-Unis.
Données~: \texttt{fm\_regmap.PERIODS}~; sources dans le registre du chapitre.}\label{fig:fm:regulators-and-licences:periods}
\end{omfigure}

Le cabinet londonien ne relève que d'un régime, un ensemble de contrôles portant sur l'interdiction générale britannique~; Chicago en ajoute trois (l'enregistrement en tant que broker-dealer, l'adhésion
à un \omterm{def:fm:regulators-and-licences:sro}{organisme d'autorégulation}, sauf si l'on n'opère que sur ses propres bourses, et la Market Access Rule)~; Hong Kong ajoute un contrôle de licence. Le délai légal de décision le plus long du parcours est
celui du Royaume-Uni, six mois pour un dossier complet et douze pour un dossier incomplet~; la décision de la SEC sur un broker-dealer intervient en 45 jours, ou jusqu'à 210 jours si
la SEC ouvre une procédure et la prolonge. Les délais légaux commencent lorsque le dossier est complet~: les mois consacrés à rédiger le programme d'activité, à recruter les personnes approuvées et à capitaliser
l'entité passent en premier, et ces délais sont la partie du parcours que l'entreprise ne maîtrise pas.

\begin{remark}[Garder la carte à jour]\label{rem:fm:regulators-and-licences:upkeep}
Chaque ligne de la carte est un fait daté portant sur une loi qui évolue. La valeur de la carte tient à sa mise à jour~: un responsable pour chaque juridiction, une date de révision pour
chaque ligne et un rapport de lignes périmées que la fonction conformité lit. Un an après le contrôle de ce chapitre, chaque ligne sera périmée.
\end{remark}

\section{Construction~: la carte réglementaire}\label{bld:fm:regulators-and-licences:regmap}

\begin{build}
\bfield{Rôle} Pas de construction au sens du projet en cours~: le sujet est le droit. L'outil d'analyse conserve la carte sous forme de données, vérifiées et interrogeables.
\bfield{Interface} \texttt{firm.regmap}~: \texttt{load}, \texttt{check}, \texttt{query}, \texttt{regimes}, \texttt{stale}, \texttt{longest\_decision},
\texttt{ALGO\_CHECKLIST}~; données dans \texttt{data/desk/regmap.csv}.
\bfield{Règles} Chaque ligne a une référence, un mois de validité et une source~; les statuts sont obligatoires~: obligation ou contrôle~; rien n'est déduit au-delà des
lignes.
\bfield{Tests d'acceptation} \texttt{code/firm/regmap/tests/}~: requêtes, régimes et délais les plus longs sur une carte de trois lignes~; le contrôle du schéma et le rapport de lignes
périmées.
\bfield{Pour aller plus loin} Lignes pour Singapour, le Japon et l'Australie avec leurs sources~; exigences de \omterm{def:fm:the-proprietary-market-making-firm:ownfunds}{fonds propres} par ligne (chapitre~2)~; une comparaison entre deux versions
de la carte.
\end{build}

\begin{omsources}
\item Directive 2014/65/UE (MiFID II), articles 2, 4, 5, 7, 17 et 34.
\item Financial Services and Markets Act 2000, sections 19 et 55V.
\item Securities Exchange Act of 1934, section 15~; 17 CFR 240.15b9-1~; Investment Advisers Act of 1940, section 203~; NFA, registration categories.
\item Securities and Futures Commission, licensing~: types of regulated activity.
\end{omsources}

\section{Exercices}

\begin{exercise}[$\star$]\label{exo:fm:regulators-and-licences:1}
Un cabinet ne négocie que pour compte propre sur des actions européennes, par l'intermédiaire d'un courtier, sans accès électronique direct ni algorithme. Déterminer s'il relève de l'exemption
de MiFID~II. Dire ce qui change s'il devient membre de la bourse.
\end{exercise}

\begin{exercise}[$\star$]\label{exo:fm:regulators-and-licences:2}
Nommer les trois caractéristiques d'une technique de trading algorithmique à haute fréquence dans MiFID~II.
\end{exercise}

\begin{exercise}[$\star$]\label{exo:fm:regulators-and-licences:3}
Quel est le délai le plus long que la SEC peut prendre pour une demande d'enregistrement en tant que broker-dealer au titre de la section 15(b)(1)~?
\end{exercise}

\begin{exercise}[$\star\star$]\label{exo:fm:regulators-and-licences:4}
Que modifie l'amendement de 2023 à la Rule 15b9-1 pour une société de trading pour compte propre qui opère sur une bourse et également hors bourse~?
\end{exercise}

\begin{exercise}[$\star\star$]\label{exo:fm:regulators-and-licences:5}
Une société américaine gère un fonds de contrats à terme pour des investisseurs extérieurs et donne des conseils de trading. Quelles catégories de la CFTC la carte désigne-t-elle, et par l'intermédiaire de qui
s'enregistre-t-elle~?
\end{exercise}

\begin{exercise}[$\star\star$]\label{exo:fm:regulators-and-licences:6}
Pourquoi la carte marque-t-elle les lignes de \omterm{def:fm:regulators-and-licences:own}{négociation pour compte propre} du Royaume-Uni et de Hong Kong comme des contrôles et non comme des obligations~?
\end{exercise}

\begin{exercise}[$\star\star\star$]\label{exo:fm:regulators-and-licences:7}
\emph{Programmation.} Interroger la carte pour un teneur de marché de l'UE utilisant une technique à haute fréquence. Combien de lignes et de régimes, et quelle est la
décision légale la plus longue~?
\end{exercise}

\begin{exercise}[$\star\star\star$]\label{exo:fm:regulators-and-licences:8}
\emph{Trouver la faille.} «~Nous ne négocions que notre propre argent, donc nous n'avons besoin d'aucune licence nulle part.~»
\end{exercise}

\section{Problème~: trois bureaux}

\begin{problem}[{Problème du week-end --- trois bureaux}]\label{pb:fm:regulators-and-licences:1}
Une société de trading pour compte propre londonienne prévoit des bureaux à Chicago et à Hong Kong et demande à son directeur des opérations le parcours réglementaire.

\medskip\noindent\textbf{Partie I --- Le périmètre.}
\begin{enumerate}
\item Définir le \omterm{def:fm:regulators-and-licences:perimeter}{périmètre réglementaire} et l'\omterm{def:fm:regulators-and-licences:perimeter}{agrément}.
\item Définir la \omterm{def:fm:regulators-and-licences:own}{négociation pour compte propre} et une \omterm{def:fm:regulators-and-licences:own}{entreprise d'investissement}.
\item Énoncer l'exemption de MiFID~II pour la \omterm{def:fm:regulators-and-licences:own}{négociation pour compte propre} et ses exceptions.
\item Définir le passeport européen (\omterm{def:fm:regulators-and-licences:passport}{passporting}) et dire ce qu'il offre à une entreprise agréée dans l'UE.
\end{enumerate}

\medskip\noindent\textbf{Partie II --- Les États-Unis.}
\begin{enumerate}[resume]
\item Définir un \omterm{def:fm:regulators-and-licences:sro}{organisme d'autorégulation} et en nommer deux.
\item Énoncer les règles d'enregistrement en tant que broker-dealer et d'adhésion, ainsi que leur calendrier.
\item Définir un \omterm{def:fm:regulators-and-licences:cftc}{swap dealer}, un \omterm{def:fm:regulators-and-licences:cftc}{commodity pool operator} (CPO) et un \omterm{def:fm:regulators-and-licences:cftc}{introducing broker}.
\item Qu'ajoute la Market Access Rule~?
\end{enumerate}

\medskip\noindent\textbf{Partie III --- La carte.}
\begin{enumerate}[resume]
\item Que porte chaque ligne de la carte, et pourquoi~?
\item Interroger les trois profils et donner lignes et régimes.
\item Donner les délais légaux de décision sur le parcours.
\item Que montre le rapport de lignes périmées un an plus tard~?
\item Que la carte ne dit-elle pas à l'entreprise~?
\end{enumerate}

\medskip\noindent\textbf{Partie IV --- Le plan.}
\begin{enumerate}[resume]
\item Quelles obligations découlent, dans l'UE, pour les algorithmes de l'entreprise~?
\item Que faut-il faire avant que le délai légal d'une demande ne commence à courir~?
\item Quelle étape du parcours l'entreprise maîtrise-t-elle, et laquelle non~?
\item Comment garder la carte à jour~?
\item Pourquoi Hong Kong est-elle un contrôle et non une obligation dans la carte~?
\item Énoncer le \emph{résultat nommé}~: les régimes que touchent les activités de l'entreprise lorsqu'elle ajoute une deuxième et une troisième juridiction, et le délai
légal de décision le plus long sur son parcours vers la négociation.
\item En deux phrases, rédiger le plan.
\end{enumerate}
\end{problem}

\section{Questions d'entretien}

\begin{interviewq}[$\star$ \iqroles{trader}]\label{iq:fm:regulators-and-licences:1}
Pourquoi une société de trading pour compte propre qui n'a aucun client a-t-elle besoin d'une licence dans l'UE~?
\end{interviewq}

\begin{interviewq}[$\star$ \iqroles{développeur}]\label{iq:fm:regulators-and-licences:2}
Que signifie juridiquement, pour votre entreprise, se livrer au trading algorithmique dans l'UE~?
\end{interviewq}

\begin{interviewq}[$\star\star$ \iqroles{risque}]\label{iq:fm:regulators-and-licences:3}
Votre entreprise veut négocier des swaps avec des banques pour son propre compte. Que vérifieriez-vous aux États-Unis~?
\end{interviewq}

\begin{interviewq}[$\star\star$ \iqroles{développeur}]\label{iq:fm:regulators-and-licences:4}
Concevez un modèle de données pour une carte réglementaire qui puisse être audité.
\end{interviewq}

\begin{interviewq}[$\star\star$ \iqroles{trader, risque}]\label{iq:fm:regulators-and-licences:5}
Qu'est-ce que le passeport européen, et que devient-il lorsqu'un pays quitte le marché unique~?
\end{interviewq}

\begin{interviewq}[$\star\star\star$ \iqroles{risque}]\label{iq:fm:regulators-and-licences:6}
Planifiez le parcours réglementaire d'une entreprise qui passe de la négociation de son propre argent à la gestion de fonds extérieurs, aux États-Unis et dans l'UE.
\end{interviewq}
